La selección tecnologica del proyecto se ha realizado atendiendo a tres criterios principales: la necesidad de construir una aplicación web full-stack mantenible, la conveniencia de reutilizar una única base de código para cliente y servidor y la necesidad de integrar servicios externos de soporte para mapas, geocodificación, imágenes y representación de QR.

\subsection{Tecnologías de desarrollo web}

La aplicación se desarrolla sobre \textbf{Next.js}, aprovechando el modelo basado en \textit{App Router}, componentes de servidor y \textit{Route Handlers} para combinar interfaz y operaciones de backend dentro del mismo proyecto \citep{nextjsdocs2026,nextjsroutehandlers2026}. Sobre esta base se utiliza \textbf{React} como librería de componentes de interfaz \citep{reactdocs2026} y \textbf{TypeScript} para tipar contratos, respuestas y servicios de dominio, reforzando la detección temprana de errores \citep{typescriptdocs2026}.

Para la capa visual se emplea \textbf{Tailwind CSS}, que permite construir interfaces responsive a partir de clases de utilidad reutilizables \citep{tailwinddocs2026}. Las transiciones y ciertos elementos de realce visual se apoyan en \textbf{Motion for React}, sucesor de Framer Motion para el ecosistema React \citep{motionreactdocs2026}.

\subsection{Persistencia y acceso a datos}

La persistencia de datos se resuelve con \textbf{PostgreSQL}, elegido por su robustez como base de datos relacional y por su adecuación para modelar usuarios, clientes, catálogo, visitas, pedidos y estados de negocio enlazados entre si \citep{postgresdocs2026}. Sobre ella se utiliza \textbf{TypeORM} como capa ORM, con entidades tipadas y migraciones versionadas que permiten evolucionar el esquema sin perder trazabilidad \citep{typeormentities2026,typeormmigrations2026}.

Este enfoque encaja especialmente bien con el proyecto, ya que la aplicación necesita integridad referencial, control de estados y reconstruccion repetible del esquema en distintos entornos.

\subsection{Autenticación y seguridad}

La autenticación del sistema se apoya en \textbf{Auth.js}, integrando un proveedor de credenciales, sesiones JWT y callbacks de control para enriquecer el contexto de usuario \citep{authjsdocs2026}. A esta capa se añaden protecciones propias del proyecto:

\begin{itemize}
    \item control de acceso por rol sobre rutas privadas;
    \item rate limiting aplicado sobre la API y sobre el flujo de login;
    \item trazabilidad de accesos y acciones administrativas;
    \item compatibilidad mínima de navegador desde el \texttt{proxy} de la aplicación.
\end{itemize}

Las contraseñas no se almacenan en claro, sino mediante verificación hash en el servidor.

\subsection{Mapas, geocodificación y rutas}

El soporte geográfico del proyecto combina varias tecnologías. La visualización y selección interactiva de ubicaciones se realiza con \textbf{Leaflet} y su adaptación a React \citep{leafletdocs2026}. La resolución de direcciones a coordenadas se apoya en \textbf{Nominatim} \citep{nominatimapi2026}, mientras que el cálculo de geometría de rutas y trayectos se apoya en \textbf{OSRM} \citep{osrmdocs2026}.

Estas integraciones no constituyen un módulo de negocio independiente, pero son esenciales para las funcionalidades ya implementadas de \texttt{M2} y \texttt{M4}, especialmente en la geocodificación de clientes, la previsualización de rutas y la estimación de reparto al cliente.

\subsection{Imágenes externas y QR}

El almacenamiento externo de imágenes se resuelve mediante \textbf{Cloudinary}, tanto para la imagen de perfil de usuario como para las imágenes del catálogo administrado por el distribuidor \citep{cloudinaryupload2026}. Para la representación del QR operativo de cada pedido se utiliza \textbf{QuickChart}, que permite generar imágenes QR a partir de una carga textual estable \citep{quickchartqr2026}.

Esta decisión evita introducir una dependencia adicional de renderizado de QR en el frontend y facilita que la misma referencia pueda reutilizarse para visualización, impresión y validación en reparto.

\subsection{Entorno de desarrollo y acceso desde dispositivos}

Durante el desarrollo local se utiliza \textbf{Docker Desktop} para levantar la base de datos y reproducir el entorno de trabajo de forma estable \citep{dockerdocs2026}. En una fase anterior del proyecto se utilizo \textbf{Cloudflare Tunnel} para exponer temporalmente la aplicación local \citep{cloudflaretunnel2026}. Sin embargo, en la iteración actual se ha pasado a usar preferentemente el sistema de \textbf{Port Forwarding} integrado en \textbf{Visual Studio Code}, sustentado sobre \textbf{Microsoft Dev Tunnels} \citep{vscodeportforwarding2026,microsoftdevtunnelsfaq2026}.

Este cambio responde a criterios prácticos de desarrollo: evita depender de una herramienta adicional externa al editor, resulta gratuito dentro del flujo habitual de trabajo, mantiene un proceso más directo para pruebas desde móvil y permite reutilizar la misma URL del túnel cuando se conserva el mismo \textit{dev tunnel}, lo que resulta especialmente útil para mantener instalada la PWA en dispositivos de prueba sin tener que reinstalarla o regenerar continuamente el punto de acceso.

\subsection{Aplicación web progresiva}

La aplicación incorpora una capa básica de \textbf{PWA}. En la implementación actual esto se concreta en:

\begin{itemize}
    \item un manifiesto web generado desde \path{app/manifest.ts};
    \item accesos directos instalables hacia panel, perfil, catálogo, pedidos y rutas mediante \path{/shortcuts/...};
    \item iconos específicos para instalación;
    \item registro de un \textit{service worker} propio;
    \item una estrategia sencilla de cache de \textit{app shell} y recuperación desde cache cuando la red falla.
\end{itemize}

Los accesos directos del manifiesto no sustituyen al control de autorización. Cada enlace se resuelve sobre rutas internas que redirigen al destino adecuado según el rol autenticado, de modo que la PWA instalada puede ofrecer accesos rápidos sin duplicar lógica de permisos ni exponer rutas privadas a usuarios no autenticados.

No se trata todavía de una PWA avanzada con sincronización offline completa, pero si de una base coherente con el uso móvil previsto para la aplicación.

La estabilidad del enlace de pruebas externas también influye en esta experiencia. Poder reutilizar una URL persistente del túnel de desarrollo simplifica la instalación continuada de la PWA en dispositivos móviles, ya que el acceso habitual puede mantenerse sobre el mismo origen durante las sesiones de prueba.

\subsection{Herramientas de documentación y apoyo}

La memoria del proyecto se mantiene en \textbf{LaTeX}, por su adecuación para documentos académicos extensos y por la facilidad con la que gestiona referencias cruzadas, tablas, figuras y bibliografía. Para diagramas y apoyo visual se han utilizado herramientas como \textbf{Draw.io} e \textbf{Inkscape}.

En conjunto, la pila tecnologica elegida permite sostener una aplicación web moderna, modular y evolutiva, manteniendo una relación directa entre requisitos, implementación y despliegue.
